% !TEX root = ../main.tex
\addcontentsline{toc}{chapter}{Conclusions}
\chapter*{Conclusions}
\markboth{Conclusions}{Conclusions}
\label{Conclusions}

In this thesis we studied the noise of PSR~J1640+2224 in the DR2low
dataset of the EPTA, with the aim of understanding the origin of the
degeneracy between the DM variations and the scattering variations that
is observed in the posterior distributions of this pulsar.

We first validated Discovery against Enterprise, and we found that the two
codes give the same posterior distributions on the same data. The
comparison between DR2low and DR2new+ showed that the low-frequency data
of LOFAR are necessary to constrain the chromatic noise, since without
them the posteriors of the DM and scattering parameters are dominated by
the priors. With DR2low the posterior of the DM parameters has two peaks,
and the contours between the DM and the scattering parameters are
L-shaped. This structure does not disappear if the chromatic index of the
scattering variations is left free, or if a band noise in the LOFAR band
is added to the model. The band noise is not detected in the real data,
and we could only set an upper limit on its amplitude.

We then used simulations to test the possible origins of this structure.
When the simulated data contain only the three processes of the model, the
injected parameters are correctly recovered and no degeneracy is found.
The frequency coverage of DR2low is therefore sufficient to distinguish
the DM variations from the scattering variations. We then injected
processes that are not included in the model. A GWB with the amplitudes
measured by the PTA collaborations produces only weak L-shaped tails in
the contours between the DM and the scattering parameters, while a
stronger GWB is absorbed by the intrinsic red noise, which becomes
stronger and steeper. A band noise in the LOFAR band biases the DM
parameters toward higher amplitudes and lower spectral indices, in the
region of the high-amplitude peak of the real data. When both processes
are injected the two effects are found together. None of these
simulations however reproduces the whole structure of the real posterior,
with its two peaks and its poorly constrained scattering parameters.

The simulations also showed that a band noise can bias the DM parameters
without being detected when it is added to the model, because of the
large uncertainties of the LOFAR TOAs. The upper limit obtained from the
real data is compatible with this case, so the presence of band noise in
the real data cannot be excluded. The origin of the degeneracy thus
remains open. A natural continuation of this work is to test other
unmodeled processes and other noise realizations, and to repeat the
analysis on other pulsars of DR2low.

These results show that an unmodeled process can bias the parameters of
all the other noise components of a pulsar, including the intrinsic red
noise, which has the same achromatic signature as the GWB. An accurate
noise model for each pulsar is therefore an essential step in the search
for the GWB.